\section{One-rate regular and stable passages}\label{sec:regular}

We use the scalar-state convention of Section~\ref{sec:additive}.
All holomorphic functions below are bounded on fixed complex neighborhoods
of the closed slow interval $[0,1]$, the state disk, and the ordinary
parameter boxes.  All evaluations use smaller neighborhoods with room.
An amplitude bound, once chosen, is fixed independently of subsequent
charge and asymptotic requests.

The proof follows the three stages of Section~\ref{sec:additive}:
actual continuation, a finite charge expansion, and dispersal of its
coefficients into packets.  Here the slow coefficient construction
uses one inverse large variable and has no surviving inner endpoint
kernel.

Let $q$ be a positive large affine combination of flag nodes and bounded
ordinary coordinates, with fixed coefficients on the large nodes.  Write
$i$ for its leading index.  As in the additive construction, convert every
participating primary basis node to a dependent node, and saturate the
flag.  Thus powers of these nodes are exact shift monomials times bounded
analytic factors.  The statement also permits the flag extensions of
Section~\ref{sec:calculus}.

\begin{theorem}[One-rate passage packets]\label{thm:regular-packets}
The following two passage outputs satisfy the packet conditions of
Definition~\ref{def:packet} and the differentiated primitive conditions of
Assumption~\ref{ass:calculus-primitive}.
\begin{enumerate}
\item A regular base flow with endpoint layers,
\begin{equation}\label{eq:regular-layer-model}
 \begin{gathered}
 z_s=V(s,z)+qW(s,E,H,z),\qquad W(s,0,0,z)=0,\\
 E_\nu=e_\nu e^{-a_\nu qs},\qquad
 H_\nu=d_\nu e^{-b_\nu q(1-s)}.
 \end{gathered}
\end{equation}
The weights are positive multiples of a fixed charge unit, and the
amplitudes are sufficiently small.  The base flow from the chosen input
disk is analytic and regular throughout $[0,1]$, with room for a small
perturbation.
\item A stable linear generator with bounded slow remainder,
\begin{equation}\label{eq:regular-stable-model}
 v_s=-qv+g(s,v),
\end{equation}
with input in a sufficiently small disk about zero.
\end{enumerate}
Ordinary parameters may occur analytically in all the displayed data.
The estimates hold for every fixed number of independent retained-free
and ordinary derivatives, and for the controlled input tests of
Assumption~\ref{ass:calculus-primitive}.  No exceptional unsafe-column
loss is required.  For each fixed finite request the estimates have
uniform constants and far-regime thresholds on the controlled families
of Definition~\ref{def:calculus-controlled-family}, independently of
the individual tests' return times to the ordinary germ.
This primitive library is retestable for elimination
in the sense of Definition~\ref{def:elimination-retestable}, with its
fixed local sector determinations.
\end{theorem}

A passage may first be reversed to put it in the stable direction in
Equation~\eqref{eq:regular-stable-model}.  The model in
Equation~\eqref{eq:regular-layer-model}
may also be applied on finitely many smaller slow intervals.  Rescaling
an interval multiplies $q$ by its fixed length and changes the endpoint
amplitudes accordingly.  A sufficiently short fixed interval has the
required regular base flow on any bounded analytic state box.

\subsection{Actual passage maps}

Write $z=\Phi_s(y)$, where $\Phi_s$ is the base flow in
Equation~\eqref{eq:regular-layer-model} and $\Phi_0=\id$.  Regularity and the
available room give an analytic inverse near the compact family under
consideration.  Consequently
\[
 y_s=q\widetilde W(s,E,H,y),\qquad
 \widetilde W=(\partial_y\Phi_s)^{-1}
 W(s,E,H,\Phi_s(y)),
\]
and $\widetilde W(s,0,0,y)=0$.  The endpoint operation $\Phi_1$ is
analytic on a larger disk than will be used.  It therefore suffices to
treat $V=0$; we again call the transformed field $W$.

On the flag sectors through band $i$, the side estimates give
\begin{equation}\label{eq:regular-leading-sector}
 \Re q\ge c\Re Z_i,\qquad \Re q\gg\log|q|,
\end{equation}
with the first inequality on the owning fringe and the corresponding
right-half-plane dominance on earlier bands.  Put $D=qs$ and use the
polygon
\begin{equation}\label{eq:regular-polygon}
 0,\quad S,\quad q-S,\quad q,\qquad S=C\log|q|.
\end{equation}
Here $C$ is a sufficiently large fixed constant.  Real parts increase,
and $D/q$ stays within $O(S/|q|)$ of $[0,1]$.  On each horizontal end the
nearby layer is integrable with bound proportional to its amplitude.
On the middle both endpoint layers have an exponential margin, giving,
for sufficiently large $C$,
\begin{equation}\label{eq:regular-layer-integral}
 \sup_\gamma(|E|+|H|)+
 \int_\gamma(|E|+|H|)\,|\dd D|
 \le C_0\max_\nu(|e_\nu|,|d_\nu|).
\end{equation}
For example, its length is $O(|q|)$ and its left-layer contribution is
at most $O(|q|e^{-a_{\min}S})$; the right layer is identical.  Opposite
layers on the end segments are controlled by
$\Re q-S\gg S$.  The vanishing of $W$ at $(E,H)=0$ gives the same
bound for the integral of its state Lipschitz constant.  The integral
equation for $y_D=W$ thus keeps the solution in the prescribed smaller
state disk, with a fixed margin, by choosing the amplitudes small.

Locally freeze $S$, vary the endpoints analytically, and solve the same
integral equation.  This proves holomorphic dependence.  Two permissible
lengths give the same solution: deform their polygons into one another
through polygons obeying the same bounds, and use analytic continuation
and uniqueness.  This continuation agrees with the real passage and is
holomorphic across the real ray.  No cutoff is involved in this actual
map.

For Equation~\eqref{eq:regular-stable-model}, on real $0\le s\le1$ the integral
equation is
\[
 v(s)=e^{-qs}v_{\mathrm{in}}+
 \int_0^s e^{-q(s-r)}g(r,v(r))\,\dd r.
\]
The forcing contribution is bounded by $\sup|g|/\Re q$ while the
solution remains in its box.  A continuation argument gives existence
and a uniform smaller disk for large $q$ and small fixed input, and
ordinary analytic ODE dependence gives holomorphy on these sectors.
These arguments apply uniformly on smaller ordinary parameter disks.

\subsection{The finite charge expansion}

We first use unit charge, so $a_\nu,b_\nu$ are positive integers and
$P=e^{-q}$. For a general charge unit, replace $P$ by the corresponding
power and change the constants below. The next lemma gives the expansion
that will be dispersed into a packet tree.

\begin{lemma}[Finite charge description]\label{lem:regular-charges}
For every integer $K\ge0$, the layer output has, on the owning
fringe, a description
\begin{equation}\label{eq:regular-charge-expansion}
 z_{\mathrm{out}}=
 \sum_{k=0}^K P^k\sum_{\ell=0}^k q^\ell
          F_{k\ell}^{[J]}(1/q)
 +O\bigl((1+|q|)^{C_K}|P|^{K+1}+e^{-c_K|q|}\bigr).
\end{equation}
Each $F_{k\ell}$ is a formal Gevrey-1 series, analytic in the ordinary
inputs. Here $[J]$ denotes truncation at a locally frozen integer
$J\asymp\delta|q|$, with $\delta>0$ sufficiently small for the fixed
request.  The stable output has the same conclusion with only
$\ell=0$.  Estimates persist under fixed ordinary derivatives on
smaller disks.  Individual constants and optimal proportions may
depend on $K$ and the fixed coefficient request.
\end{lemma}

To obtain this expansion, separate the layer monomials that remain
constant along the passage from those that decay towards one endpoint.
For layer multiindices put
\[
 A=a\cdot m,\qquad B=b\cdot n,\qquad \Delta_{mn}=B-A.
\]
Along the actual profiles,
\[
 E^mH^n=e^m d^n P^B e^{(B-A)qs}.
\]
The terms with $A=B$ are therefore constant in $s$ apart from their
slow coefficients. Their accumulated drift supplies the powers of $q$
in Equation~\eqref{eq:regular-charge-expansion}. The other terms can be
removed successively using the nonzero charge gap. We implement this
separation by the layer derivation
\[
 \mathcal D_{\rm lay}=-\sum a_\nu E_\nu\partial_{E_\nu}
                  +\sum b_\nu H_\nu\partial_{H_\nu},
\]
which has eigenvalue $\Delta_{mn}$ on $E^mH^n$.
Let $\Pi$ denote its zero-eigenvalue projection. For the stable model,
a formal change about its slow graph leaves a linear fiber equation;
its endpoint damping supplies the factor $P$ without an accumulated
balanced drift. This will give the same charge expansion with $\ell=0$.

\subsubsection{Formal normalization}

\begin{lemma}[Slow-strip formal normalization]\label{lem:regular-formal}
Put $h=1/q$.
For the layer model there are unique formal series in $h$ and the fast
variables of the form
\[
 y=u+t(s,E,H,u;h),\qquad \Pi t=0,\qquad
 u_D=\mathcal C(s,E,H,u;h),\qquad \Pi\mathcal C=\mathcal C,
\]
determined by
\begin{equation}\label{eq:regular-layer-normalization}
 \mathcal C=\Pi W(s,E,H,u+t),\qquad
 (h\partial_s+\mathcal D_{\rm lay})t=(1-\Pi)
       \bigl(W(s,E,H,u+t)-t_u\mathcal C\bigr).
\end{equation}
There is no constant balanced coefficient.  For each fixed left
multiindex the full right tail converges in a common fixed fast radius,
and conversely with left and right interchanged.  On any fixed finite
collection of such tails and balanced coefficients, the coefficient of
$h^j$ has norm at most $C^{j+1}j!$.  The pure transformations and their
state derivatives are small on smaller fixed fast disks.

For the stable model there are formal series
\[
 T(s,u;h)=v_0(s;h)+u+\sum_{m\ge2}t_m(s;h)u^m,
 \qquad b(s;h)=g_v(s,v_0(s;h)),
\]
with a convergent common $u$-disk norm and the same Gevrey bound,
satisfying
\begin{align}
 h\partial_sv_0&=-v_0+hg(s,v_0),\label{eq:regular-graph}\\
 [h\partial_s+(m-1)(-1+hb)]t_m
  &=h[g(s,T)-g(s,v_0)-b(T-v_0)]_{u^m},\quad m\ge2.
 \label{eq:regular-fiber}
\end{align}
Here $v_0$ and all $t_m$ start at positive $h$-order.  All these
statements allow any fixed ordinary and state derivatives on smaller
disks.  Constants may depend on the finite request; the fast amplitude
and input bounds need not shrink with that request.
\end{lemma}

\begin{proof}
Differentiating $y=u+t$ gives
\[
 y_D=\mathcal C+h t_s+\mathcal D_{\rm lay}t+t_u\mathcal C.
\]
Multiplication by a balanced monomial preserves imbalance, so
$\Pi(t_u\mathcal C)=0$.  This proves
Equation~\eqref{eq:regular-layer-normalization}, including its signs.  Every
monomial of $W$ contains a layer factor.  The equations are therefore
triangular in total fast degree, even at $h$-order zero; nonzero
imbalances have the fixed gap $|\Delta_{mn}|\ge1$.  The coefficient with
$m=n=0$ vanishes and stays zero under the recursion.

Here are the norm details needed in the presence of infinite tails.
Use the sum of coefficient sup norms times fixed fast radii, with sup
norms on slightly larger slow and state domains.  For a pure side the
drift is absent.  Its composition operator has small Lipschitz norm
because $W$ and $W_y$ have a fast factor.  Charge division has norm at
most one, so a small fixed fast radius gives its analytic majorant and
small state derivative on a smaller state disk.

Fix a nonzero left multiindex $m$ and induct on its total degree.  The
only dependence of $W$ on the current-left tail is linear convolution
with
\[
 W_y(s,0,H,u+t_{0,*}),
\]
which has no constant $H$ coefficient.  At current left order, a
differentiated factor in $t_u\mathcal C$ has lower left order, except
that the known pure derivative $t_{0,*,u}$ may multiply a current
balanced coefficient.  That balanced coefficient is determined on the
finite list $B=A$, in increasing right degree, before it can occur in
the higher right tail.  Thus no derivative of the unknown current
tail occurs.  First solve the finite list $B\le N$.  For $N$ large
relative to $A$, the remaining diagonals satisfy
$|\Delta_{mn}|\ge c(1+|n|)$.  The inverted convolution on this tail has
norm at most $C/N$, hence is absorbable.  Its forcing consists of
convergent lower-left tails and the finite initial list.  Fix one fast
radius strictly inside the pure radius throughout; increasing the
threshold, rather than decreasing that radius, absorbs every such
convolution.  The argument with the sides interchanged is identical.
State derivative margins can be allotted summably by total one-side
degree, so the common remaining state disk is nonempty.

To control the formal slow derivatives, work modulo $h^{J+1}$, with
weighted norm
\[
 \sum_{j=0}^J (\varepsilon/J)^j\|F_j\|_j \quad (J\ge1),
\]
where the slow domains decrease linearly with $j$ across one fixed
allotted margin.  Product norms are submultiplicative.  Cauchy's
estimate across one step of the slow margin bounds $\partial_s$ by
$C J$, so $h\partial_s$ has norm $O(\varepsilon)$.  Choose
$\varepsilon$ small after the finite coefficient request has been
fixed.  It is then absorbed after each nonzero charge division.  The
pure small equation and the current-left high-tail convolution have
the same bounds uniformly in $J$.  Lower-order state derivatives use
the already allotted margins, and the balanced coefficients are
obtained algebraically.  This gives a bound independent of $J$ for
each requested weighted norm.  Taking $J=j$ gives
$C(j/\varepsilon)^j\le C_1^{j+1}j!$.  Margins for the slow strip,
like the state margins, may be allotted summably over successive
fixed-side degrees.  This proves all the asserted layer bounds
without asserting convergence when both fast indices grow without
restriction.

For the stable model, the desired conjugacy equation is
\begin{equation}\label{eq:regular-stable-conjugacy}
 hT_s+(-1+hb)uT_u=-T+hg(s,T).
\end{equation}
Its constant coefficient is Equation~\eqref{eq:regular-graph}, its linear
coefficient is the definition of $b$, and its higher coefficients
give Equation~\eqref{eq:regular-fiber}.  The graph equation is a small formal
fixed-point equation in the same weighted norm: $h\partial_s$ is
small and $h g$ has small analytic Lipschitz norm.  After the graph is
known, divide the fiber equations by $m-1$.  This diagonal inverse is
bounded on the absolute $u$-coefficient norm for $m\ge2$; it cancels
the factor $m-1$ in the $hb$ term.  Taylor composition about $v_0$ is
bounded and Lipschitz on a fixed sufficiently small $u$ disk, with
range strictly inside the field disk.  The nonlinear remainder has
no constant or linear fiber term.  Thus the small $h$ multiplier and
the small weighted slow derivative give a convergent majorant for
all fiber degrees at once.  At $h=0$ the graph and nonlinear fiber
terms vanish.  The weighted estimate again gives the factorial
bound.  Cauchy estimates on remaining parameter and state margins
give the stated fixed derivatives.
\end{proof}

\subsubsection{Optimal truncation and comparison with the actual flow}

For a fixed finite request choose a sufficiently small constant
$\delta>0$ and truncate every series needed for it at a common order
$J\asymp\delta|q|$.  Locally this integer is frozen.  The weighted
proof above applies at a dummy $h$ radius $c/J$ larger than the used
$|h|$ by a fixed factor.  All operations in the proof are bounded on
that dummy disk.  Formal cancellation through degree $J$ and Cauchy's
estimate therefore bound the substitution residual by
$O(e^{-c|q|})$ on smaller domains.  A layer diagonal factor can be
absorbed by one final fixed shrink of the fast radius, since
$|\Delta_{mn}|$ grows at most linearly in the fast degree.  This observation
also controls fixed derivatives of the residual.  A smaller
proportional cutoff, or two overlapping fixed-factor choices, changes
each retained expression by $O(e^{-c'|q|})$.

\begin{proof}[Proof of Lemma~\ref{lem:regular-charges}]
For the layer model, retain the normalization modulo the monomial
ideal generated by $E^mH^n$ with $A,B\ge K+1$.  This consists of
finitely many fixed-one-side tails, so Lemma~\ref{lem:regular-formal}
applies.  Write $t=t_L+t_R+t_{\mathrm{mix}}$ in this quotient.
Along the polygon, a mixed monomial obeys
\begin{equation}\label{eq:regular-mixed-smallness}
 |e^{-AD-B(q-D)}|\le |P|^{\min(A,B)},
 \qquad 0\le\Re D\le\Re q.
\end{equation}
Hence $t_{\mathrm{mix}}=O_K(P)$ in modulus, whereas $t_L+t_R$
is uniformly small by the fixed amplitude choice.  Taylor-expand
$W(s,E,H,u+t_L+t_R+t_{\mathrm{mix}})$ in its last mixed increment
through degree $K$.  Its actual remainder is $O_K(|P|^{K+1})$.
This finite Taylor calculation is legitimate on the larger slow and
state domains; it does not require $t_{\mathrm{mix}}$ to be small on
an unrestricted fast bidisk.

All its low-charge formal coefficients cancel in the normalization
equation, giving residual $O(e^{-c_K|q|})$ after optimal truncation.
Any remaining analytic fast sum whose monomials have $A,B\ge K+1$
has actual value $O_K(|P|^{K+1})$.  To see this directly, choose from
each such monomial componentwise subindices with left and right
charge at least $K+1$ and minimal for that property.  Their charges
are bounded by $K+1$ plus the largest weight, so only finitely many
subindices occur.  Extracting them gives $|P|^{K+1}$ by
Equation~\eqref{eq:regular-mixed-smallness}; the remaining absolute series
is bounded by the fixed fast norm.  This also applies to the finite
mixed Taylor operations just used.

The retained balanced drift is exactly
\[
 \mathcal C(s,E,H,u;h)
       =\sum_{k=1}^K P^k C_k^{[J]}(s,u;h)
\]
on the actual layer trajectory, since every nonconstant balanced
monomial has $A=B=k\ge1$.  Invert the input transformation and
solve this approximate drift.  The input inverse is regular: its
pure part is close to identity and its mixed part tends to zero.
The drift moves $u$ by $O_K(|q||P|)$, keeping it in its disk.
Along the actual polygon the Lipschitz integral for $W$ is bounded
by Equation~\eqref{eq:regular-layer-integral}.  Comparing the transformed
approximate solution with the true solution by the integral
inequality therefore costs at most a constant times the polygon
length $O(|q|)$.  Polynomial losses are absorbed by reducing the
exponential constant or by the displayed power in
Equation~\eqref{eq:regular-charge-expansion}.

At the two endpoints the fast arguments are respectively
\[
 (E,H)=(e,(d_\nu P^{b_\nu})_\nu),\qquad
 (E,H)=((e_\nu P^{a_\nu})_\nu,d).
\]
Treat $P$ temporarily as independent.  On
$|P|\le c_K/(1+|q|)$ the endpoint inverses and the slow drift
$u_s=q\sum_{k=1}^KP^kC_k^{[J]}(s,u;h)$ remain in fixed bounded
boxes.  Their output is bounded and analytic there.  Cauchy's
estimate gives a Taylor remainder bounded by
$C_K(1+|q|)^{K+1}|P|^{K+1}$ on the actual exponentially smaller
$P$.  Each coefficient of degree $k$ uses finitely many endpoint
inversion jets and at most $k$ nested drift integrals, because every
drift insertion has positive $P$ degree.  Changing $\dd D$ to
$q\,\dd s$ places its $q$ factor outside the coefficient integral.
Induction on $k$ thus gives a polynomial in $q$ of degree at most
$k$, with coefficients formed by bounded analytic operations and
integrals over fixed subintervals of $[0,1]$.  The final $\Phi_1$
evaluation has the same property.  The Gevrey composition and
regular-inverse estimates of Lemma~\ref{lem:additive-gevrey-operations},
or its one-variable dummy-disk proof, apply to these finite operations.
Integration on a fixed interval preserves the same coefficient
bound.  This proves Equation~\eqref{eq:regular-charge-expansion}.  In
particular no large $q$ occurs inside an order-zero inverse: the
entire drift has positive charge.

For the stable model, take optimal polynomials $T^{[J]},b^{[J]}$
from Lemma~\ref{lem:regular-formal}, using a common sufficiently
small proportion.  The transformation is $\id+O(h)$ and has a
regular inverse on the fixed smaller input disk.  Set
\[
 u(s)=\exp\left(-qs+\int_0^s b^{[J]}(r;h)\,\dd r\right)
                  (T^{[J]}(0,\cdot;h))^{-1}(v_{\mathrm{in}}).
\]
Since $b^{[J]}$ is bounded, damping and a fixed input margin keep
$u(s)$ in the fiber disk.  The residual of
$v_{\mathrm{app}}(s)=T^{[J]}(s,u(s);h)$ in the $D$ equation is
$R=O(e^{-c|q|})$.  If $L$ bounds $|g_v|$ on the working disk,
variation of constants and the scalar integral inequality give
\[
 |v(1)-v_{\mathrm{app}}(1)|
 \le \frac{|q|}{\Re q-L}\sup_{0\le s\le1}|R(s)|.
\]
The logarithm of this loss is $O(\log|q|)$ in
Equation~\eqref{eq:regular-leading-sector}, hence it is absorbed in the
exponential error.  The approximate endpoint is
\begin{equation}\label{eq:regular-stable-endpoint}
 T^{[J]}\left(1,
 P\exp\left(\int_0^1b^{[J]}(s;h)\,\dd s\right)
       (T^{[J]}(0,\cdot;h))^{-1}(v_{\mathrm{in}});h\right).
\end{equation}
Taylor expansion in $P$ takes place on a fixed disk and uses only
bounded operations with regular inverses.  The same Gevrey rules
give the desired coefficients, with no outside power of $q$.
Ordinary derivative estimates in both constructions follow on
smaller fixed parameter disks from these uniform analytic estimates.
\end{proof}

\subsection{Dispersal and independent derivatives}

We complete the proof of Theorem~\ref{thm:regular-packets}, making
the retirement of numerical truncations explicit.  Write
\begin{equation}\label{eq:regular-rate-split}
 q=q_0(1+\Theta),\qquad q_0=A_0Z_i,\qquad A_0>0,
 \qquad h_0=q_0^{-1}.
\end{equation}
By the dependent-node assumptions, $\Theta$ is a finite sum of
strict-small shift monomials with first index greater than $i$,
times bounded analytic factors.  The exact formula for $f_i$
uses only nodes of index at least $\sigma(i)>i$ and the bounded
backgrounds.  Thus $h_0$ is itself a strict-small shift monomial
with first index $\sigma(i)$.

Up to band $i$ use the actual holomorphic output and zero prefixes
on preceding transitions.  At the $i$ fringe take the charge
description to any sufficiently large fixed $K$.  Since
$\Re q\ge c\Re Z_i$ and $\log|q|=o(\Re Z_i)$, its charge
remainder pays every prescribed fixed depth by increasing $K$.
Its optimal error pays every such depth too: on that fringe
$\Re Z_i/|Z_i|\to0$, while $|q|\asymp|Z_i|$.  Keep the affine
exponential $P^k=e^{-kq}$ exactly as its flag shift, including its
bounded factor.  In particular changing a ratio argument in a later
Taylor operation does not change this already extracted shift.

For a coefficient in Equation~\eqref{eq:regular-charge-expansion}, factor
$q^\ell=q_0^\ell(1+\Theta)^\ell$.  The first factor is an
explicit shift translation.  Absorb the unit into the bounded
coefficient and rewrite
\[
 \sum_{a\ge0}F_aq^{-a}
       =\sum_{a\ge0}F_a(1+\Theta)^{-a}h_0^a.
\]
On one fixed small complex disk of the ratio arguments this changes
the Gevrey bound by at most a factor $C^a$.  Freeze the optimal
order using a fixed-factor representative of $|q_0|$, independently
of those ratio arguments.  It follows that the coefficient and all
its fixed ratio derivatives have a uniformly bounded optimal
description on smaller ratio and ordinary disks.

Proceed through the later flag indices in order.  At index $j$,
Taylor-expand precisely those analytic small arguments whose first
index is $j$, to a sufficiently large fixed degree for the requested
budget.  Set these arguments to zero in subsequent coefficient
evaluations, retaining every other argument.  Their remainders pay
arbitrary $j$-depth because each current-first monomial costs at
least $e^{-c\Re Z_j}$ on its owning fringe.  Until $j=\sigma(i)$,
keep the $h_0$ series optimally evaluated.  Its numerical log formula
uses only retained inputs throughout those bands.  At
$j=\sigma(i)$ keep only a fixed sufficiently large $h_0$ degree:
the Gevrey estimate and the small optimal proportion give
\[
 \left|\sum_{a=J}^{N_{\mathrm{opt}}}F_a(\Theta)h_0^a\right|
       \le C_J|h_0|^J
\]
for each fixed $J$.  This pays arbitrary fixed current depth.
Thereafter the coefficient uses finite formal degrees only; its
numerical optimal order has disappeared.  Every shift is split at
the current index and its remaining tail is passed to the next
band.  This is the single-variable instance of
Proposition~\ref{prop:additive-dispersal}, with no inner kernels.

These recursive operations also verify the support condition.
Charges at $i$ are nonnegative integers; at each fixed charge there
are finitely many explicit translations.  The further supports are
Taylor powers of a finite collection of lexicographically strict
positive shifts and the powers of $h_0$, each dispersed at its
first index.  Only finitely many current-first degrees contribute
below a given ceiling.  Repeating this observation at a fixed
prefix proves iterated left-finiteness; coincident exponent tuples
are added with finite multiplicity.  At the last level the
coefficients are analytic ordinary germs.  No simultaneous
convergence of the full tree has been used.

For completeness consider cutoffs, raw growth, and derivatives.
Choose the representative moduli for $|q_0|=A_0e^{\Re f_i}$ in
overlapping fixed-factor bins.  On each bin use a frozen integer
cutoff and hence a holomorphic expression in the retained inputs.
Neighboring expressions differ, with each fixed derivative on
smaller domains, by
\begin{equation}\label{eq:regular-cutoff-cost}
 O\bigl(e^{-c e^{\Re f_i}}\bigr).
\end{equation}
Smoothly interpolate on overlaps as in the packet definition.
These sources occur only in bands
$i<j\le\sigma(i)$, where $f_i$ is retained in its exact
triangular formula; their cost $U=e^{\Re f_i}$ is an allowed
simple source cost.  The flag cost estimates give
$\Re Z_j+f_j=o(U)$ there.  Thus all fixed shift normalizations,
raw factors, and interpolation derivative losses have logarithmic
size $o(U)$.  The actual expression on band 1 has no such source.

With explicit shift factors removed, a fixed coefficient is bounded
on smaller analytic argument disks.  Remaining shifts have modulus
at most $\exp(C\Re Z_j)$ on band $j$ by side dominance.  This
proves raw growth on the whole band with no $D_{j-1}$ term.  In
particular every admissible column is safe, including all the
entering collars required by Definition~\ref{def:packet}.

To differentiate in independent coordinates, freeze local cutoff
choices and take complex neighborhoods of radius
$\exp(-C f_j(u))$ in the retained free and ordinary inputs.
The finite triangular formulas bound the nodes of index at least
$j$ and their fixed sensitivities by $\exp(C'f_j(u))$.
An earlier active logarithm such as $f_i$ is evaluated by its exact
linear and bounded-background formula and has the same sensitivity
bound.  Its exponential $q_0$ need not have that absolute bound:
keeping the change in $f_i$ bounded instead keeps $q_0$ within a
fixed factor and preserves the frozen optimal-cutoff margin.
By increasing $C$, and beginning on a slightly wider polynomial
buffer, these neighborhoods preserve the slow, state, small-ratio,
and angular margins used above.  Cauchy's estimate costs only
$\exp(C_r f_j(u))$ for each fixed derivative order $r$.
Since $\Re Z_j\gg f_j(u)^N$ for every fixed $N$ up to the side,
deeper initial transition budgets absorb these losses.  At a
cutoff source they are absorbed by $U$ in
Equation~\eqref{eq:regular-cutoff-cost}; derivatives of the interpolating
partitions have the same permitted logarithmic loss.  Applying
the same argument on both sides of each transition proves the
differentiated transitions and all fixed differentiated anti-CR
defect bounds.

\subsection{Controlled inputs and uniformity}

\begin{lemma}\label{lem:regular-controlled-inputs}
For each fixed finite request, the preceding primitive estimates are
uniform on a controlled family of
Definition~\ref{def:calculus-controlled-family}.  Finite real-Taylor
agreement supplies the geometric comparisons for evaluation of the
independent-variable estimates.  Under pullback, differentiated source
bounds hold for holomorphic tests and for the smooth chart sheets with
the differentiated source bounds supplied by
Theorem~\ref{thm:calculus}.
\end{lemma}

\begin{proof}
Fix the request, the flag and rate recipe, and the common controlled
family data.  In particular fix a compact ordinary box inside the
larger analytic domains, with a common positive margin, on the full
forward domain of the tests.  The field sup norms, base-flow inverse
margin, state and fast disks, and all ordinary Cauchy constants can
be fixed on these domains.  The polygon integral and the stable
variation-of-constants bound use only these norms, the fixed weights,
and the displayed sector inequalities.  Hence their constants and
required lower bounds on the large rates are common to the family.

For the formal estimates, the pure-side contraction constants and
the high-tail convolution bound use those same common analytic
norms.  For each of the finitely many one-side indices required,
choose its finite initial list and high-tail threshold once.  Choose
the weighted slow-derivative constant $\varepsilon$ and the reserved
slow and state margins for this finite collection once as well.
The stable divisors $m-1$ and the nonzero layer charge gaps are
fixed.  Thus the Gevrey constants, optimal proportions, residual
bounds, endpoint inverse margins, and finite charge remainder
constants are uniform over the ordinary box.  These are estimates
of analytic functions at the current ordinary arguments; none
waits for a test to return closer to its limiting ordinary point.

The common finite real-jet bounds and real-Taylor error constants,
with the common small-jet modulus where a slow estimate requires it,
give the same flag comparisons and angular room used in the dispersal
proof.  Its ratio disks, side constants, and Cauchy-loss exponents
can consequently be chosen for the whole family.  After the finite
request has fixed all powers to absorb, the inequalities between
the flag scales give common far-regime thresholds for those
absorptions, expressed in the finitely many numerical rate and gap
inequalities in the request.  This uses the controlled family's comparisons on the
full forward domain, rather than eventual entry of each test into
a smaller ordinary neighborhood.  The same argument applies to
cutoff overlaps and their fixed derivatives.  The fixed-factor bins
have bounded overlap, so the number of local source terms is uniform
as well.  For a coefficient
already constructed, retain its chosen convention on a larger
argument disk; any smaller comparison cutoff used in a new finite
request is reconciled with it by
Equation~\eqref{eq:regular-cutoff-cost}.  No uniform choice over
all requests is asserted, nor is a uniform rate of return of the
ordinary values to their germ point needed.

For antiholomorphic derivatives, use the source estimates of the frozen
holomorphic formulas and their explicit cutoff overlaps. Finite real-Taylor
agreement preserves their domains; it does not supply a source estimate
for an arbitrary smooth input. A holomorphic ordinary test contributes
no input defect. A smooth chart sheet supplied by
Theorem~\ref{thm:calculus} already has the required differentiated
source bounds.

The remaining factors in the pullback chain rule are the primitive raw
and Cauchy factors bounded above and the controlled sheet derivatives.
At an intrinsic cutoff source their logarithmic losses are $o(U)$ by
Equation~\eqref{eq:regular-cutoff-cost} and the comparisons above.
For an inherited sheet source, use its permitted source cost and the
comparisons of Definition~\ref{def:packet} and
Equation~\eqref{eq:calculus-source-fringe}; no restriction to a simple
cost is imposed on that source. At a level-$p$ join the later-rate and Cauchy losses are $o(V_p)$;
a fixed current normalization is retained as $C\Re Z_p$. Thus
Lemma~\ref{lem:calculus-source-pullback} applies, with a transition
budget larger than the fixed current losses when the defect comes from
a transition, or a sufficiently vertical collar for an independently
established $e^{-cV_p}$ defect. Band~1 uses holomorphic primitives and
holomorphic chart sheets and has no source. This proves the source
statement with the same controlled-family uniformity.
\end{proof}

After any first-index extraction all formulas use only the retained
free inputs and ordinary variables, and the fixed local determinations
agree whenever those inputs agree.

\subsection{Retesting the coefficient germs}

We verify the additional assertion of
Definition~\ref{def:elimination-retestable}.  The formal
normalizations in Lemma~\ref{lem:regular-formal} are uniquely determined
by the fixed analytic fields and weights: imbalance division and
balanced projection determine the layer coefficients, while the graph
and fiber recursions determine the stable coefficients.  No numerical
large input or integration anchor is an initial datum in these formal
recursions.  At fixed charge and formal degree,
Lemma~\ref{lem:regular-charges} then uses only their fixed coefficients,
analytic endpoint inversion at the same regular root, and integration
over the fixed slow interval.  The resulting ordinary analytic germs
are consequently independent of the numerical anchor along a flag
test.

More explicitly, perturb the retained old free inputs by arbitrarily
small absolute amounts while preserving the affine flag relations,
regime gaps, limiting ordinary data, and chosen sector determinations.
The exact affine shifts and the ratio formulas in
Equation~\eqref{eq:regular-rate-split} retain the same coefficient
expressions.  Every finite Taylor extraction therefore uses the same
formal coefficient germ.  The numerical optimal order may change,
but on overlaps its effect has the precision in
Equation~\eqref{eq:regular-cutoff-cost}; at $\sigma(i)$ it is removed
altogether, leaving only fixed formal degrees.  It thus cannot appear
in a terminal coefficient.  There are no additional anchor-dependent
kernels in this construction.  Refining the flag by unused nodes
inserts zero-prefix identity transitions; necessary exact log
substitutions merely rewrite the same monomials before pullback.
On any saturated extension the same finite-request proofs above
apply to that rewritten affine rate and its retained log formulas.
Theorem~\ref{thm:calculus} supplies the stated calculus on those
extensions, and the coefficient argument just given supplies its
terminal-germ invariance.  This concerns only the selected local
lifted determinations.  It proves retestability and completes
Theorem~\ref{thm:regular-packets}.
